% Three ways of making a computer larger. Drawn in place of the stacked
% PowerPoint slices, which shared one anchor and could not be edited.
\begin{tikzpicture}[remember picture, overlay, x=1bp, y=-1bp,
  shift={(current page.north west)}]
\path[fill=cFFFFFF] (0,0) rectangle (960,540);
\path[fill=cEE7628,even odd rule] (57.18,55.73) ..controls(57.96,55.73)and(58.6,56.37)..(58.6,57.15) --(58.6,85.5) ..controls(58.6,86.28)and(57.96,86.91)..(57.18,86.91) ..controls(56.4,86.91)and(55.76,86.28)..(55.76,85.5) --(55.76,57.15) ..controls(55.76,56.37)and(56.4,55.73)..(57.18,55.73);
\node[anchor=north west, inner sep=0pt] at (56.69,477.21) {\includegraphics[width=73.42bp, height=37.42bp]{figures/fundamentals-of-hpc/p03-1.png}};

\fill[cFCE4D4, rounded corners=10bp] (88,128) rectangle (900,228);
\fill[cEA7600, rounded corners=8bp] (104,146) rectangle (400,210);
\node[anchor=center, text=white] at (252,178)
  {\fcal\bfseries\fontsize{16bp}{19bp}\selectfont One big node};
\node[anchor=north west, text=c404040, align=left] at (420,152)
  {\fcal\fontsize{15bp}{20bp}\selectfont
   Many CPUs, cores, or accelerators in one machine.\\
   One operating system, one memory.};

\fill[cCAE6F9, rounded corners=10bp] (88,244) rectangle (900,344);
\fill[c325F8B, rounded corners=8bp] (104,262) rectangle (400,326);
\node[anchor=center, text=white] at (252,294)
  {\fcal\bfseries\fontsize{15bp}{18bp}\selectfont Several nodes, shared memory};
\node[anchor=north west, text=c404040, align=left] at (420,268)
  {\fcal\fontsize{15bp}{20bp}\selectfont
   Separate operating systems.\\
   Software makes the memory look shared.};

\fill[cCDDFD1, rounded corners=10bp] (88,360) rectangle (900,456);
\fill[c2CA055, rounded corners=8bp] (104,378) rectangle (400,442);
\node[anchor=center, text=white] at (252,410)
  {\fcal\bfseries\fontsize{16bp}{19bp}\selectfont Many nodes};
\node[anchor=north west, text=c404040, align=left] at (420,376)
  {\fcal\fontsize{15bp}{20bp}\selectfont
   Each node has its own system and memory.\\
   The nodes exchange messages over a network.};
\end{tikzpicture}
